\documentclass[11pt]{article}
\usepackage[letterpaper,margin=0.85in,headheight=14pt]{geometry}
\usepackage[T1]{fontenc}
\usepackage{lmodern}
\usepackage{microtype}
\usepackage{amsmath,amssymb,amsthm,mathtools}
\usepackage{booktabs,longtable,tabularx,array}
\usepackage{enumitem}
\usepackage{xcolor}
\usepackage{fancyhdr}
\usepackage[hidelinks]{hyperref}
\usepackage{listings}
\usepackage{tcolorbox}
\tcbuselibrary{breakable,skins}

\definecolor{navy}{RGB}{28,52,91}
\definecolor{pale}{RGB}{239,244,251}
\definecolor{green}{RGB}{32,92,65}
\definecolor{red}{RGB}{128,38,38}
\definecolor{graybox}{RGB}{247,247,247}

\hypersetup{
  pdftitle={Automated Logical Deciders Research Notebook, Entry 001},
  pdfauthor={Brian Tenneson},
  pdfsubject={Living research notebook on certificate-indexed multiagent logical decision},
  pdfkeywords={ALD, ATP, independence, multiagent, lemma pool, certificates, machine learning}
}

\newtheorem{theorem}{Theorem}[section]
\newtheorem{proposition}[theorem]{Proposition}
\newtheorem{lemma}[theorem]{Lemma}
\newtheorem{corollary}[theorem]{Corollary}
\theoremstyle{definition}
\newtheorem{definition}[theorem]{Definition}
\newtheorem{protocol}[theorem]{Protocol}
\newtheorem{question}[theorem]{Open Question}
\theoremstyle{remark}
\newtheorem{remark}[theorem]{Remark}
\newtheorem{conjecture}[theorem]{Conjecture}

\newcommand{\T}{\mathcal T}
\newcommand{\M}{\mathcal M}
\newcommand{\Pool}{\mathcal L}
\newcommand{\MetaPool}{\mathcal K}
\newcommand{\Ev}{\mathcal E}
\newcommand{\N}{\mathbb N}
\newcommand{\Prov}{\mathsf{PROVED}}
\newcommand{\Refut}{\mathsf{REFUTED}}
\newcommand{\Ind}{\mathsf{INDEPENDENT}}
\newcommand{\Inc}{\mathsf{INCONSISTENT}}
\newcommand{\Both}{\mathsf{BOTH}}
\newcommand{\Unknown}{\mathsf{UNKNOWN}}
\newcommand{\Ver}{\operatorname{Ver}}
\newcommand{\Con}{\operatorname{Con}}
\newcommand{\Supp}{\operatorname{supp}}
\newcommand{\ALD}{\mathsf{ALD}}

\pagestyle{fancy}
\fancyhf{}
\lhead{ALD Research Notebook}
\rhead{Entry 001 -- August 2, 2026}
\cfoot{\thepage}
\renewcommand{\headrulewidth}{0.4pt}

\setlist[itemize]{leftmargin=1.5em}
\setlist[enumerate]{leftmargin=2em}
\setlength{\parskip}{0.38em}
\setlength{\parindent}{0pt}

\newtcolorbox{notebox}[1][]{
  breakable,
  colback=pale,
  colframe=navy,
  boxrule=0.6pt,
  arc=2pt,
  left=8pt,right=8pt,top=6pt,bottom=6pt,
  title=#1,
  fonttitle=\bfseries
}
\newtcolorbox{decisionbox}[1][]{
  breakable,
  colback=graybox,
  colframe=green,
  boxrule=0.6pt,
  arc=2pt,
  left=8pt,right=8pt,top=6pt,bottom=6pt,
  title=#1,
  fonttitle=\bfseries
}

\begin{document}

\begin{titlepage}
\centering
\vspace*{0.55in}
{\Huge\bfseries Automated Logical Deciders\par}
\vspace{0.12in}
{\LARGE\bfseries Research Notebook\par}
\vspace{0.18in}
{\Large Entry 001: Certificate-Indexed Halting, Independence Detection, and a Shared Lemma Pool\par}
\vspace{0.55in}
{\Large Brian Tenneson\par}
\vspace{0.08in}
{\large M.S., Applied Data Science\par}
{\large Currently Unaffiliated\par}
\vfill
\begin{minipage}{0.86\textwidth}
\small
\textbf{Notebook status.} This is a living research record, not a polished replacement for the full ALD paper. It records the motivating insight, working definitions, theorem sketches, literature findings, originality boundaries, frozen design decisions, experiment plans, and unresolved questions as of the opening entry. Future entries should preserve prior entries and append dated changes rather than silently rewriting the research history.
\end{minipage}
\vfill
{\large Started August 2, 2026, 9:55 AM Pacific\par}
\vspace{0.18in}
{\small Copyright \textcopyright\ 2026 Brian Tenneson. All Rights Reserved.\par}
\end{titlepage}

\tableofcontents
\clearpage

\section{Notebook purpose and provenance}

This notebook begins after the observation that an Automated Logical Decider should not be modeled merely as two ATPs, one searching for a proof of $c$ and one searching for a proof of $\neg c$. If both searches fail to halt, the computer has not detected independence. It has only failed to find either proof. The research question is therefore:

\begin{notebox}[Opening research question]
What finite, independently checkable evidence can license a multiagent system to report that a conjecture is proved, refuted, inconsistent with the base theory, independent of the base theory, or still unknown?
\end{notebox}

The second central observation is that the agents should not operate as isolated portfolios. Every agent should be able to add verified lemmas to a shared pool and use every compatible verified lemma at every round. This makes cooperation part of the mathematical object, not merely an implementation detail.

\subsection{Connection to Version 47}
Version 47 already distinguishes staged proof search, theorem birth, halting, finite consequence closure, preserved invariants, proof certificates, and bounded unknown outcomes. The present project reorganizes those ingredients around a richer terminal-status theory. The intended extension is not ``an ATP that tries harder.'' It is a certificate-indexed multiagent system whose different outcomes require different kinds of evidence.

\section{Entry 001: the motivating obstruction}

Fix an object theory $\T$, a conjecture $c$, and a declared metatheory $\M$ capable of checking claims about $\T$. Consider five roles:
\begin{enumerate}
  \item a positive agent $A_{+}$ searching for $\T\vdash c$;
  \item a negative agent $A_{-}$ searching for $\T\vdash\neg c$;
  \item a contradiction agent $A_{\bot}$ searching for $\T\vdash\bot$ or, more specifically, for derivations of both $c$ and $\neg c$;
  \item an independence agent $A_{I}$ searching in $\M$ for certificates of $\T\nvdash c$ and $\T\nvdash\neg c$;
  \item optional specialist agents for model finding, finite saturation, invariant discovery, relative consistency, proof translation, and certificate minimization.
\end{enumerate}

The original obstruction is that $A_{+}$ and $A_{-}$ may both run forever. Their silence has no unique logical interpretation. The replacement is to make nonderivability a positive search target in the metatheory.

\begin{decisionbox}[Frozen conceptual distinction]
Program termination, exhaustive search closure, and logical settlement are different. A timeout may stop a program; finite saturation may certify closure of a declared search space; only a status-appropriate certificate may settle a logical claim.
\end{decisionbox}

\section{Working formal architecture}

\subsection{Typed evidence}
\begin{definition}[Evidence types]
For an ALD instance $(\T,c,\M)$, the basic evidence types are:
\begin{align*}
E_{+} &: \T\vdash c,\\
E_{-} &: \T\vdash\neg c,\\
E_{\bot} &: \T\vdash\bot,\\
E_{nc} &: \M\vdash (\T\nvdash c),\\
E_{n\neg c} &: \M\vdash (\T\nvdash\neg c).
\end{align*}
An evidence ledger $\Ev_r$ records all verifier-accepted evidence available after round $r$.
\end{definition}

\begin{definition}[Status map]
A conservative status map is defined from the ledger by:
\begin{align*}
\Prov &\quad\text{if }E_{+}\in\Ev,\ E_{-},E_{\bot}\notin\Ev,\\
\Refut &\quad\text{if }E_{-}\in\Ev,\ E_{+},E_{\bot}\notin\Ev,\\
\Both &\quad\text{if }E_{+},E_{-}\in\Ev,\\
\Inc &\quad\text{if }E_{\bot}\in\Ev,\\
\Ind &\quad\text{if }E_{nc},E_{n\neg c}\in\Ev\text{ and no accepted inconsistency evidence defeats the declared interpretation},\\
\Unknown &\quad\text{otherwise.}
\end{align*}
\end{definition}

The ledger is more informative than a single label. In particular, if an inconsistent theory yields both $c$ and $\neg c$, the system should preserve both proofs and the inconsistency evidence rather than overwrite one result with the other.

\subsection{The universally shared lemma pool}
\begin{definition}[Provenance-aware lemma record]
An object-level lemma record is a tuple
\[
\lambda=(\ell,\Delta,\pi,o,r),
\]
where $\ell$ is a formula, $\Delta$ is an explicit finite support set of assumptions, $\pi$ is a verifier-accepted certificate of $\T\cup\Delta\vdash\ell$, $o$ records the originating agent, and $r$ is the first publication round. A metatheoretic record is defined analogously in $\M$.
\end{definition}

\begin{definition}[Round pool]
Let $\Pool_r$ be the object-level pool and $\MetaPool_r$ the metatheoretic pool after round $r$. Every agent reads a snapshot of all compatible records in $\Pool_r\cup\MetaPool_r$ and may submit proposed additions during round $r+1$. Only independently verified proposals are published.
\end{definition}

\begin{protocol}[Round-synchronous ALD]
At each round $r$:
\begin{enumerate}
  \item freeze a read-only snapshot $(\Pool_r,\MetaPool_r,\Ev_r)$;
  \item allow every agent to perform one declared unit of search using every compatible record in the snapshot;
  \item collect proposed lemmas and status certificates;
  \item independently verify every proposal in a deterministic audit order;
  \item publish accepted additions to form $(\Pool_{r+1},\MetaPool_{r+1},\Ev_{r+1})$;
  \item apply the frozen halting policy to the new evidence ledger;
  \item checkpoint the private agent states and global audit log.
\end{enumerate}
\end{protocol}

\begin{remark}
The one-round delay is not logically necessary, but it makes causality, reproducibility, and race-condition auditing much cleaner. An asynchronous implementation can later be related to the synchronous model by recording a serializable publication order.
\end{remark}

\section{Initial theorem log}

The following are working theorems for the ALD theory. Several are already formalized more fully in the accompanying paper; the notebook records why they matter and what remains to be sharpened.

\begin{theorem}[Pool monotonicity]
For all $r$, $\Pool_r\subseteq\Pool_{r+1}$, $\MetaPool_r\subseteq\MetaPool_{r+1}$, and $\Ev_r\subseteq\Ev_{r+1}$.
\end{theorem}
\begin{proof}
The protocol publishes only verified additions and never deletes an accepted record. Hence each structure is monotone by construction.
\end{proof}

\begin{theorem}[Object-pool soundness]
If every published object record $(\ell,\Delta,\pi,o,r)$ is accepted only when the object verifier confirms $\T\cup\Delta\vdash\ell$, then every use of that record in a branch whose assumptions include $\Delta$ is sound relative to $\T$ and that branch.
\end{theorem}
\begin{proof}
The verifier establishes the recorded sequent. Weakening permits its use under any assumption context containing $\Delta$. No other branch may treat $\ell$ as a theorem of $\T$ unless the support is discharged.
\end{proof}

\begin{corollary}[Safe universal visibility]
Every agent may see every pool record without compromising soundness, provided admissibility checks prevent an agent from using a record outside a context containing its recorded support.
\end{corollary}

\begin{theorem}[Support-discharge promotion]
Suppose $(\ell,\Delta,\pi,o,r)\in\Pool_r$ and for every $\delta\in\Delta$ the pool contains a verified proof of $\T\vdash\delta$. Then a verifier can construct a certificate of $\T\vdash\ell$, and $\ell$ may be promoted to a base-theory lemma with empty support.
\end{theorem}
\begin{proof}
Substitute the verified derivations of the support assumptions into the certificate $\pi$ and check the resulting composed derivation.
\end{proof}

\begin{theorem}[Status soundness]
Assume each status is emitted only after its designated sound verifier accepts the required evidence. Then every emitted conclusive status is correct relative to the declared object theory, metatheory, and status semantics.
\end{theorem}
\begin{proof}
By cases on the status. For $\Prov$, $\Refut$, and $\Inc$, the accepted object certificate supplies the corresponding derivation. For $\Ind$, the metatheoretic verifier accepts both nonderivability claims. The status map adds no stronger claim than the accepted evidence supports.
\end{proof}

\begin{theorem}[Unknown safety]
Resource exhaustion, heuristic stagnation, agent crash, or failure to locate a certificate may produce only $\Unknown$ unless some conclusive evidence has already been accepted.
\end{theorem}
\begin{proof}
These operational events are not proof objects for any conclusive logical claim. The frozen status policy therefore maps them to $\Unknown$.
\end{proof}

\begin{theorem}[Positive independence detection]
If the ALD accepts certificates $\pi_c$ and $\pi_{\neg c}$ such that
\[
\Ver_{\M}(\pi_c,\ulcorner\T\nvdash c\urcorner)=1
\quad\text{and}\quad
\Ver_{\M}(\pi_{\neg c},\ulcorner\T\nvdash\neg c\urcorner)=1,
\]
then the ALD may finitely emit $\Ind$ without waiting for the positive and negative object-level searches to diverge forever.
\end{theorem}
\begin{proof}
The two accepted certificates are positive finite evidence for both nonderivabilities. The independence status follows by its definition.
\end{proof}

\begin{theorem}[Fair certificate discovery]
Fix a recursively enumerable class of finite certificates for each conclusive evidence type. Suppose every valid finite certificate is eventually proposed by at least one agent, every proposal is eventually checked, and the verifier terminates on every finite proposal. If a conclusive certificate exists in one of the covered classes, then the ALD eventually accepts it and reaches the corresponding conclusive ledger state.
\end{theorem}
\begin{proof}
The valid certificate appears after finitely many proposals by fairness. It is then checked after finitely many additional rounds and accepted. Ledger monotonicity preserves it thereafter.
\end{proof}

\begin{theorem}[No universal total ALD]
For any sufficiently expressive recursively axiomatized theory with undecidable consequence problem, there is no computable sound ALD that returns a correct conclusive status for every sentence.
\end{theorem}
\begin{proof}[Proof idea]
A total sound ALD deciding $\Prov$ versus non-$\Prov$ for every sentence would decide theoremhood, contradicting undecidability. Adding more agents and certificate classes can enlarge finite coverage but cannot remove the computability barrier.
\end{proof}

\section{How independence can be detected now}

The crucial move is to search for finite witnesses of nonderivability rather than interpreting nontermination. Several certificate families can support this.

\subsection{Paired models}
If the logic is sound and complete and the ALD produces:
\[
\mathfrak M_{+}\models \T\cup\{c\},
\qquad
\mathfrak M_{-}\models \T\cup\{\neg c\},
\]
then neither $c$ nor $\neg c$ follows from $\T$. This is the semantic pattern represented by the TPTP SZS status \emph{NoConsequence}: one model supports the conjecture with the axioms and another supports its negation with the axioms.

\subsection{Dual finite saturation}
When a complete calculus on a finite effective domain reaches verified saturation for both $\T\cup\{c\}$ and $\T\cup\{\neg c\}$ without contradiction, those saturated states can serve as finite evidence of satisfiability, hence of nonderivability of the opposite sides under completeness assumptions.

\subsection{Dual invariants}
An invariant may exclude every derivation of $c$ from $\T$, while another excludes every derivation of $\neg c$. The certificate must include proof that the initial state satisfies the invariant, every inference preserves it, and the excluded target violates it.

\subsection{Relative consistency or formal metaproof}
For deep independence results, the independence agent may search for a checked metaproof establishing
\[
\Con(\T)\Rightarrow\Con(\T+c)
\quad\text{and}\quad
\Con(\T)\Rightarrow\Con(\T+\neg c),
\]
or the appropriate stronger statement. The formal Lean proof of the independence of CH demonstrates that forcing-based independence reasoning can itself be kernel-checked, although automating the discovery of such arguments remains far beyond ordinary ATP search.

\begin{notebox}[Approximately 100-word solution statement]
The old method did not detect independence: two ATPs searched for $c$ and $\neg c$, and if neither proof existed both could run forever. The ALD replaces silence with positive evidence. A dedicated independence agent searches in a declared metatheory for two finite, checkable certificates: one proving that $c$ is not derivable from $\T$, and one proving that $\neg c$ is not derivable from $\T$. These certificates may be paired models, exhaustive finite saturations, dual invariants, relative-consistency arguments, or a checked formal metaproof such as forcing. Once both certificates verify, the ALD halts with $\Ind$. If they are not found, the result remains $\Unknown$.
\end{notebox}

\section{Why the shared lemma pool may change the problem}

The pool permits cross-status cooperation. A lemma discovered while trying to prove $c$ may simplify the search for $\neg c$, expose a contradiction, help construct a model, or become a premise in a metatheoretic nonderivability proof. Conversely, a model-finding agent may identify a structural invariant that prunes both proof searches. This creates several research hypotheses.

\begin{conjecture}[Cross-status reuse advantage]
There exist natural benchmark families for which a universally shared verified lemma pool yields asymptotically fewer settlement expansions than every architecture using the same agents with role-isolated pools.
\end{conjecture}

\begin{conjecture}[Independence-agent acceleration]
On some independent instances, lemmas generated by the positive and negative proof agents substantially reduce the cost of constructing the two nonderivability certificates, even though neither object-level proof search can succeed.
\end{conjecture}

\begin{conjecture}[Contradiction-agent value]
A contradiction agent that consumes the shared pool can detect hidden inconsistency earlier than a portfolio that merely waits to see whether both $c$ and $\neg c$ are independently proved.
\end{conjecture}

\subsection{A necessary warning}
Unlabeled sharing is unsound. If the positive agent temporarily assumes $c$ and proves $\ell$, the negative agent may not treat $\ell$ as a theorem of $\T$. The pool must preserve support and logical level. This provenance requirement is therefore load-bearing, not administrative.

\section{Later machine-learning guidance}

Machine learning should be introduced only after the symbolic protocol and verifiers are frozen. The learned component may rank:
\begin{itemize}
  \item which agent receives the next expansion budget;
  \item which open goal or branch is expanded;
  \item which assertion or inference rule is tried;
  \item which pool lemmas are retrieved for a given agent state;
  \item which proposed lemmas are likely to have cross-agent utility;
  \item which certificate family appears promising for the current instance.
\end{itemize}

\begin{theorem}[Guidance-neutral soundness]
Replacing a hand-designed search policy by a learned policy does not change the soundness of accepted ALD statuses, provided the learned component cannot bypass the object verifier, metatheoretic verifier, support checker, or status map.
\end{theorem}
\begin{proof}
The learned policy changes only proposal order and resource allocation. Every accepted lemma and status certificate still passes the same fixed verification conditions.
\end{proof}

\begin{remark}[Fairness protection]
Pure learned ranking may starve a necessary low-ranked certificate path forever. A practical ALD should reserve a nonzero exploration schedule, such as round-robin expansions, stochastic exploration, or an exhaustive dovetailing component, so that finite-certificate completeness is not destroyed by the learned policy.
\end{remark}

\begin{conjecture}[Cross-agent utility learning]
Training on first-use events across distinct agents yields a more transferable lemma-utility signal than training only on membership in final proofs, because it measures whether a lemma changes another search process rather than merely whether it appears in one completed certificate.
\end{conjecture}

\section{Originality research log}

\subsection{Search conducted for Entry 001}
The opening literature review checked the following neighboring areas:
\begin{enumerate}
  \item result ontologies and positive statuses for nonconsequence;
  \item distributed and multiagent theorem proving;
  \item blackboard architectures with globally visible proof states;
  \item proof and model certificates;
  \item formalized independence proofs;
  \item reusable lemma pools and lemma mining;
  \item machine-learned inference guidance.
\end{enumerate}

\subsection{Closest prior art found}
\begin{longtable}{@{}p{0.25\textwidth}p{0.66\textwidth}@{}}
\toprule
\textbf{Neighboring idea} & \textbf{What the literature already contains} \\
\midrule
Positive nonconsequence status & The TPTP SZS ontology defines \emph{NoConsequence}: some models of the axioms satisfy the conjecture and some satisfy its negation. It explicitly identifies a pair of models or saturations as output evidence. \\
Multiagent blackboard proving & Leo-III was designed with reasoning agents cooperating through a shared blackboard at term, clause, and search levels. Earlier work describes globally shared data visible to all agents. \\
Distributed clause exchange & Clause-Diffusion and Aquarius use concurrent deductive processes that partition search and exchange clauses; modified variants address fairness and proof reconstruction. \\
Lemma pools and ML & Large-theory ATP research mines millions of intermediate lemmas, reuses them across later problems, and combines them with learned relevance filtering. \\
Learned inference ranking & ENIGMA learns clause-selection guidance inside saturation-based theorem provers while the symbolic prover retains inference correctness. \\
Kernel-checked independence & The independence of CH has been formalized in Lean using Boolean-valued models and forcing for both directions. \\
\bottomrule
\end{longtable}

\subsection{Provisional novelty boundary}
The individual ingredients are not new. The strongest currently defensible originality claim is narrower:

\begin{notebox}[Candidate-original synthesis]
The reviewed sources did not reveal an established theory combining all of the following in one formal object: status-indexed object and metatheoretic agents for proof, negation-proof, contradiction, and paired nonderivability; a monotone evidence ledger separating stopping, exhaustion, and settlement; and a universally readable and writable lemma pool whose entries are independently verified, typed by logical level, and labeled by assumption support at every round.
\end{notebox}

This is not proof of worldwide novelty. A broader systematic review, citation chaining, expert consultation, and peer review are still required. In particular, multiagent blackboard systems come close on shared state, and the SZS ontology comes close on the semantic status taxonomy. The likely contribution is the \emph{certificate-indexed SIC-style synthesis and theorem package}, not the invention of cooperation or model-pair independence evidence.

\section{Experiment notebook: proposed first studies}

\subsection{Experiment A: pool ablation}
Run the same frozen agents on the same theorem set under:
\begin{enumerate}
  \item isolated private lemma pools;
  \item a shared pool without ML;
  \item a shared pool with provenance-aware learned retrieval.
\end{enumerate}
Measure settlement rate, expansions, wall time, memory, accepted lemmas, cross-agent first uses, and certificate type.

\subsection{Experiment B: status ablation}
Compare:
\begin{enumerate}
  \item proof-only ATP;
  \item dual proof/negation ATP;
  \item proof/negation/contradiction ALD;
  \item full ALD including nonderivability certificate agents.
\end{enumerate}
The main question is not merely how many proofs are found, but how many instances receive a logically justified finite status rather than $\Unknown$.

\subsection{Experiment C: synthetic independence benchmark}
Construct finite formal systems where independence is known by design and where paired saturation or invariant certificates are short. This permits testing of ALD mechanics before attempting forcing-level mathematics.

\subsection{Experiment D: Predator integration}
Treat Predator variants as proposal agents under one common verifier. Begin with the locked 50-theorem benchmark for theorem-proving performance, then build a separate ALD benchmark containing theorem, negation-theorem, inconsistent, finite-independent, and bounded-unknown instances. The theorem benchmark and ALD-status benchmark should not be conflated.

\section{Frozen decisions as of Entry 001}

\begin{decisionbox}[D1: certificate requirement]
No conclusive status is emitted solely because an agent stopped, timed out, exhausted its assigned local budget, or produced no new preferred actions.
\end{decisionbox}

\begin{decisionbox}[D2: universal pool access]
Every agent may contribute to and use the shared verified lemma pool at every round, subject to logical-level and assumption-support compatibility.
\end{decisionbox}

\begin{decisionbox}[D3: verifier separation]
Search agents, including learned agents, do not determine correctness. Independent object and metatheoretic verifiers determine acceptance.
\end{decisionbox}

\begin{decisionbox}[D4: independence means paired nonderivability]
The ALD reports independence only after accepting both $\T\nvdash c$ and $\T\nvdash\neg c$ relative to a declared metatheory and stated assumptions.
\end{decisionbox}

\begin{decisionbox}[D5: unknown remains a legitimate result]
Undecidability and certificate scarcity imply that some runs must remain $\Unknown$. The goal is to convert some previously silent cases into finite certified settlements, not to promise a universal decider.
\end{decisionbox}

\section{Open questions for subsequent entries}

\begin{question}[Best status algebra]
Should $\Inc$ dominate all other statuses, or should the ALD always output the full evidence vector $(E_{+},E_{-},E_{\bot},E_{nc},E_{n\neg c})$ together with a derived human-readable label?
\end{question}

\begin{question}[Minimal support]
Can support minimization be performed efficiently enough that the pool records near-minimal assumptions, preventing unnecessary contamination of cross-agent reuse?
\end{question}

\begin{question}[Metatheory hierarchy]
Should the independence agent itself be an ALD operating in a stronger metatheory, producing a hierarchy $\T_0\subseteq\T_1\subseteq\cdots$ of certificate-checking levels?
\end{question}

\begin{question}[Cross-agent fairness]
What is the weakest scheduler condition ensuring that no finite certificate path is permanently starved while still permitting aggressive learned allocation?
\end{question}

\begin{question}[Originality]
Which earlier systems combine model finding, proof search, and a shared blackboard with typed status evidence? A systematic citation search should focus on ``cooperative decision procedures,'' ``proof and model portfolios,'' ``blackboard theorem proving,'' ``non-theorem certificates,'' and ``independence automation.''
\end{question}

\begin{question}[Benchmark design]
How should the first ALD benchmark balance object-theory proofs, refutations, contradictions, finite independence, semantic independence, and genuinely unresolved cases without leaking target-specific certificates into training?
\end{question}

\section{Next-entry template}

Future notebook entries should append:
\begin{enumerate}
  \item date, time, and trigger;
  \item new definitions or changes to old definitions;
  \item theorem statements with proof status: proved, proof sketch, computationally tested, conjectured, or refuted;
  \item literature searches and exact novelty boundaries;
  \item experiments, frozen code/data hashes, and outcomes;
  \item decisions that are now locked;
  \item mistakes, failed approaches, and unresolved contradictions;
  \item next actions.
\end{enumerate}

\clearpage
\section*{References consulted for Entry 001}
\addcontentsline{toc}{section}{References consulted for Entry 001}
\begin{thebibliography}{99}

\bibitem{SZS}
G. Sutcliffe,
``The SZS Ontologies and Presentation Standards,''
TPTP documentation, especially the \emph{NoConsequence}, \emph{Unknown}, \emph{Timeout}, proof, model, and saturation statuses,
accessed August 2, 2026.

\bibitem{BonacinaHsiang1995}
M. P. Bonacina and J. Hsiang,
``The Clause-Diffusion Methodology for Distributed Deduction,''
\emph{Fundamenta Informaticae}, 24(1--2):177--207, 1995.
DOI: 10.3233/FI-1995-24128.

\bibitem{BonacinaReconstruction1996}
M. P. Bonacina,
``On the Reconstruction of Proofs in Distributed Theorem Proving: A Modified Clause-Diffusion Method,''
\emph{Journal of Symbolic Computation}, 21(4--6):507--522, 1996.
DOI: 10.1006/jsco.1996.0028.

\bibitem{WisniewskiBenzmuller2016}
M. Wisniewski and C. Benzm\"uller,
``Is It Reasonable to Employ Agents in Automated Theorem Proving?''
in \emph{Proceedings of ICAART 2016}, vol. 1, pp. 281--286.
DOI: 10.5220/0005824702810286.

\bibitem{KaliszykUrban2015}
C. Kaliszyk and J. Urban,
``Learning-Assisted Theorem Proving with Millions of Lemmas,''
\emph{Journal of Symbolic Computation}, 69:109--128, 2015.
DOI: 10.1016/j.jsc.2014.09.032.

\bibitem{KaliszykUrbanVyskocil2015}
C. Kaliszyk, J. Urban, and J. Vysko\v{c}il,
``Lemmatization for Stronger Reasoning in Large Theories,''
in \emph{FroCoS 2015}, LNAI 9322, pp. 341--356.
DOI: 10.1007/978-3-319-24246-0\_21.

\bibitem{ENIGMA}
J. Jakub\r{u}v and J. Urban,
``ENIGMA: Efficient Learning-Based Inference Guiding Machine,''
\emph{arXiv:1701.06532}, 2017.

\bibitem{HanVanDoorn2021}
J. M. Han and F. van Doorn,
``A Formal Proof of the Independence of the Continuum Hypothesis,''
\emph{arXiv:2102.02901}, 2021.

\end{thebibliography}

\end{document}
